\section{Metrics Pipeline and Long-Term Storage}

Metrics start close to where things happen. Kubernetes components, nodes, workloads, network
devices, storage systems, and applications all produce time-series data. The platform collects this
through a combination of exporters and monitors, some built into the stack, some added for specific
systems like SNMP devices or facility services.

Prometheus handles the collection and short-term storage. It scrapes targets, evaluates alert
rules, and makes metrics available for immediate incident response
\cite{lsst-it-k8s-cookbook,prometheus-operator,prometheus-alerting}. The configuration is adapted
for Rubin, with the right labels, selectors, and alerting choices, but the underlying tooling comes
from the Prometheus ecosystem, which is widely used and well understood.

For anything beyond the short incident-response window, Mimir provides long-term storage. It is a
horizontally scalable backend that speaks the same query language as Prometheus, so the experience
for operators and dashboard authors does not change \cite{grafana-mimir}. Under the hood, it uses
object storage and runs as a distributed service with separate components for ingestion, querying,
compaction, and caching.

This gives the team two distinct capabilities. Recent metrics support fast triage---you can see
what is happening right now and what changed in the last hour. Older metrics support a different
kind of question: how does tonight compare to last month? Is this cluster behaving differently from
the others? Is this a new problem or a slow drift that has been building for weeks? That longer
view is especially important during commissioning, when the baseline is still being established and
slow changes are easy to miss.
